\chapter{The Squeeze}
% OUTLINE Ch4 "The Squeeze (numerics' face: the convergence structure)".
% Laws: INTUITION LAW (labeled blocks, adds-no-claim), CITATION
% (Archimedes' Measurement of a Circle at Cavendish grade; Cauchy 1821;
% Bolzano 1817/Weierstrass arithmetization; Richardson 1911/1927 mesh
% extrapolation), LENS EXHIBIT 2 OF 2 rides here (Beastmaster's
% arrival-read): Vol 1's count-that-stops as the shape of an honest
% squeeze HALTED -- with the calibration adapted honestly: the device's
% BRACKETS are carried (proved order, documented in Vol 1); what stays
% fenced is the LIMIT reading (the device stops, it does not pass to the
% limit). Dose closes at 2. Gate: Kodo.

\begin{intuition}
Two hands closing on a fly. You never touch it --- the fly is faster
than the hands --- and yet there is a moment when you know where it is,
because the gap that must contain it has shrunk to almost nothing. The
feeling to carry into this chapter: \emph{caught} means \emph{bracketed},
never \emph{touched}, and an entire face of mathematics stands on the
difference.
\end{intuition}

The ur-instance is the oldest calculation in this series' ledgers, and
it deserves its story. Archimedes, wanting the ratio of a circle's
boundary to its width, did not compute it --- no finite arithmetic
lands it --- and did not estimate it either, in the loose sense of that
word. He trapped it. Inscribe a hexagon: its perimeter falls short of
the circle's, and honestly so, being made of chords. Circumscribe a
hexagon: too long, honestly, being made of tangents. Now double ---
twelve sides, twenty-four, forty-eight, ninety-six --- each doubling
computable by nothing but square roots, each pair of polygons a pair of
sworn statements: the circle lies between us. At ninety-six sides he
laid down the pen and published the two fences standing, in the
\emph{Measurement of a Circle}: the ratio exceeds three and ten
seventy-firsts, and falls short of three and one seventh. The number
itself was never touched, never named, never needed. What was published
was a bracket and the method that tightens it at will --- and that
publication, twenty-two centuries old, is the complete pattern of the
face this chapter teaches.

What the pattern needed, and waited two millennia for, was the license
to speak of the quarry at all. Cauchy, in the \emph{Cours d'analyse} of
1821, gave convergence its modern grammar --- including the criterion
that lets one certify a sequence closes in on \emph{something} without
ever producing the something --- and the arithmetization that followed,
from Bolzano's early rigor of 1817 through Weierstrass's lectures, made
the backbone explicit: bounded monotone sequences converge; nested
closed intervals with shrinking widths share exactly one point. On that
backbone the squeeze becomes a proof structure with a signature as
clean as the first two faces'. To prove a statement about a number you
cannot exhibit, build two sequences you can: one below the quarry, one
above, both provably closing. Every property that holds for all the
brackets passes to the prisoner. The quantifier lives in the
construction --- \emph{for every} width, a bracket narrower --- and the
conclusion is a fact about a point no step of the proof ever touched.

The face has its own variation, and naming it earns the chapter's place
in this volume. The first face varied the candidate; the second varied
the landscape; the numerical face varies the \emph{mesh}. Coarsen or
refine the grid on which a quantity is computed, and watch the readings
organize: Lewis Fry Richardson, in 1911, computing stresses in masonry
by finite differences, saw that the error of a mesh reading falls in a
lawful way with the mesh's fineness --- and that two readings at two
meshes therefore contain more than either alone, the lawful part of the
error cancelling between them, the limit approached as if deferred ---
a method he and Gaunt named outright in 1927. The numericist's wiggle
is the mesh; the analogue of the vanishing first variation is the
stabilizing of readings under refinement; and the honest output is,
always, a bracket with its width stated.

\begin{intuition}
The surveyor's version of the feeling: you never stand on the point.
You drive a stake on either side and swear only to the between --- and
your oath gets better not by touching the point but by driving the
stakes closer, each advance a computation you can actually perform.
\end{intuition}

\begin{fence}
An illustration, not an identity --- and this volume's second device
appearance, closing its dose. The first volume of this series documents
an instrument that walks a bracket toward its quarry and \emph{stops}
--- at a count its own ground earned --- publishing the two walls in
proved order and displaying, never asserting, the point between them.
That practice is offered here as the shape of this chapter's structure
worn honestly by a machine: brackets constructed, walls certified,
the point left untouched. One calibration differs from this volume's
first exhibit and is stated plainly: the instrument's brackets are its
own carried structure --- the walls and their order are checked
mechanically, as its documentation records --- so what is illustrated
is not a resemblance but a practice. What stays behind the fence is the
\emph{limit}: this chapter's theorems pass to a point at the end of an
infinite closing, and the instrument, by its own law, never does --- it
stops at its earned count and asserts the interval. Whether its stopped
bracket could ever be earned as this chapter's convergence is a question
the series leaves exactly where it leaves all such questions: owed, not
crossed. Illustration --- and here, honest practice --- but not
identity.
\end{fence}

Take stock, because the volume's tour of faces is complete. To prove
best: vary the candidate, one arbitrary wiggle carrying every rival ---
the extremal structure. To prove must-exist: read the invariant the
space imposes --- the obstruction structure. To prove exactly: find the
arrangement in which the content cancels to zero --- the identity
structure. To prove about the unreachable: trap it between computable
fences that close --- the convergence structure. Four faces, four
signatures, every one cited, every picture labeled. What remains is the
suspicion the reader has surely formed by now --- that these are not
four subjects but four lights on one --- and the next chapter takes one
proof the reader already owns and turns it under all four, to see.
